% concatenated from Med_Lattice_Linfty_260301.tex
\documentclass[a4paper,11pt]{article}
\usepackage[utf8]{inputenc}
\usepackage{a4wide}
\usepackage{latexsym,amsfonts,amsmath,amssymb,mathrsfs,url,amsthm}
\usepackage{mathtools}
\usepackage{color,graphicx}
\usepackage{lipsum}
\usepackage{hyperref}
\usepackage{xcolor}
\usepackage{cleveref}
\usepackage{bm}
\usepackage{comment}

\newtheorem{theorem}{Theorem}[section]
\newtheorem{lemma}[theorem]{Lemma}
\newtheorem{example}[theorem]{Example}
\newtheorem{assumption}[theorem]{Assumption}
\newtheorem{algorithm}[theorem]{Algorithm}
\newtheorem{remark}[theorem]{Remark}
\newtheorem{proposition}[theorem]{Proposition}
\newtheorem{definition}[theorem]{Definition}
\newtheorem{corollary}[theorem]{Corollary}
% Keywords command
\providecommand{\keywords}[1]
{
  \noindent \small	
  \textbf{Keywords:} #1
}
% AMScode command
\providecommand{\amscode}[1]
{
  \noindent \small	
  \textbf{AMS subject classifications:} #1
}
\newcommand{\rd}{\, \mathrm{d}}
\newcommand{\rand}{\mathrm{rand}}
\newcommand{\wor}{\mathrm{wor}}
\newcommand{\rms}{\mathrm{rms}}
\newcommand{\app}{\mathrm{app}}
\newcommand{\bszero}{\boldsymbol{0}}
\newcommand{\bsh}{\boldsymbol{h}}
\newcommand{\bsg}{\boldsymbol{g}}
\newcommand{\bsk}{\boldsymbol{k}}
\newcommand{\bsl}{\boldsymbol{\ell}}
\newcommand{\bsL}{\boldsymbol{L}}
\newcommand{\bsx}{\boldsymbol{x}}
\newcommand{\bsy}{\boldsymbol{y}}
\newcommand{\bsz}{\boldsymbol{z}}
\newcommand{\bsgamma}{\boldsymbol{\gamma}}
\newcommand{\bsDelta}{\boldsymbol{\Delta}}
\newcommand{\EE}{\mathbb{E}}
\newcommand{\II}{\mathbb{I}}
\newcommand{\NN}{\mathbb{N}}
\newcommand{\PP}{\mathbb{P}}
\newcommand{\RR}{\mathbb{R}}
\newcommand{\ZZ}{\mathbb{Z}}
\newcommand{\Acal}{\mathcal{A}}
\newcommand{\Hcal}{\mathcal{H}}
\newcommand{\Pcal}{\mathcal{P}}
\newcommand{\Zcal}{\mathcal{Z}}
\newcommand{\fraku}{\mathfrak{u}}
\newcommand{\Mod}[1]{\ (\mathrm{mod}\ #1)}
\newcommand{\tmod}[1]{{\;(\mathrm{mod}\; #1)}}

\newcommand{\e}{\mathbb{E}}
\newcommand{\bsone}{\boldsymbol{1}}

\DeclareMathOperator{\supp}{supp}
\DeclareMathOperator{\proj}{proj}

\newcommand{\icomp}{\mathtt{i}}
\newcommand{\abs}[1]{\left\vert#1\right\vert}
\newcommand{\norm}[1]{\left\Vert#1\right\Vert}
\newcommand{\floor}[1]{\left\lfloor #1 \right\rfloor} % floor
\newcommand{\ceil}[1]{\left\lceil #1 \right\rceil}    % ceil






\allowdisplaybreaks

%\usepackage[notref,notcite]{showkeys}
\mathtoolsset{showonlyrefs=true}

\newcommand{\peter}[1]{\begingroup\color{violet}#1\endgroup}
\newcommand{\zexin}[1]{\begingroup\color{blue}#1\endgroup}
\newcommand{\tg}[1]{\begingroup\color{teal}#1\endgroup} 
\newcommand{\josef}[1]{\begingroup\color{orange}#1\endgroup}
\newcommand{\mc}[1]{\begingroup\color{magenta}#1\endgroup}
\newcommand{\TODO}[1]{{\color{red} TODO: #1}}



\title{Worst-case $L_p$-approximation of periodic functions using median lattice algorithms}
\author{Zexin Pan\thanks{Institute of Fundamental and Transdisciplinary Research, Zhejiang University, 866 Yuhangtang Road, Xihu District, Hangzhou, Zhejiang Province, 310058, China (\url{zep002@zju.edu.cn})}, Mou Cai\thanks{Graduate School of Engineering, The University of Tokyo, 7-3-1 Hongo, Bunkyo-ku, Tokyo 113-8656, Japan (\url{caimoumou@g.ecc.u-tokyo.ac.jp}; \url{goda@frcer.t.u-tokyo.ac.jp}).}, Josef Dick\thanks{nice affiliation}, Takashi Goda\footnotemark[2], Peter Kritzer\thanks{RICAM, Austrian Academy of Sciences, Altenbergerstr. 69, 4040 Linz, Austria (\url{peter.kritzer@oeaw.ac.at})}}
\date{\today}

\begin{document}

\maketitle

\sloppy

\centerline{{\it Dedicated to Henryk Wo\'zniakowski on the occasion of his 80th birthday.}}

\begin{abstract}
We study the worst-case approximation of multivariate periodic functions from the weighted Korobov space $H_{d,\alpha,\gamma}$ with smoothness $\alpha>1/2$ in the Lebesgue norm $L_p([0,1]^d)$ for $1\le p\le\infty$. We analyze a \emph{median lattice algorithm} that reconstructs a truncated Fourier series by approximating the coefficients on a hyperbolic-cross-type index set using $R$ rank-1 lattice sampling rules with independent randomly chosen generating vectors, and then aggregating the resulting coefficient estimators via the componentwise median. For an odd number of repetitions $R>1$ and an odd prime lattice size $N$, we prove high-probability error bounds in both $L_\infty$ and $L_2$. Interpolation then yields the result for all $1 \le p\le\infty$. In particular, with a high probability, the algorithm satisfies
\[
\mathrm{err}(H_{d,\alpha,\gamma},L_p,A)\ \le\ C_{d,\alpha,\beta,\bsgamma,p}\, N^{- \alpha + (\frac12 - \frac1p)_+ + \beta },
\qquad 1 \le p\le\infty,\ \beta>0,
\]
where $(x)_+ = \max\{x, 0\}$, $N$ is the number of function evaluations, and the weights $\bsgamma$ and the constant $C_{d,\alpha,\beta,\bsgamma,p}$ are independent of $N$. For $p=\infty$, $C_{d,\alpha,\beta,\bsgamma,\infty}$ is dimension-independent under the summability condition $\sum_{j=1}^\infty \gamma_j^{1/(2\alpha)}<\infty$. These results extend recent analyses of median-based lattice approximation in $L_2$ and complement related multiple-shift lattice approaches, showing that median aggregation yields nearly optimal $L_p$-approximation rates (up to logarithmic factors and an arbitrarily small loss) in weighted Korobov spaces.
\end{abstract}



\keywords{multivariate approximation, lattice rule, weighted Korobov space}

\amscode{42A10, 65D15, 65T40}
%%%%%%%%%%%%%%%%%%%%%%%%%%%%%%%%%%

%%%%%%%%%%%%%%%%%%%%%%%%%%%%%%%%%%%%%%%%%%


\section{Introduction}

Rank-1 lattice rules are a classical and powerful tool in quasi-Monte Carlo (QMC) methods for high-dimensional problems. They have been developed extensively for numerical integration in periodic function spaces, with a mature theory covering existence, construction, randomization, and tractability; see, e.g., \cite{DKP22,DKS13,N92,SJ94}. In particular, for weighted Korobov spaces, component-by-component (CBC) constructions and their fast variants provide lattice rules that achieve optimal or near-optimal convergence rates under natural summability assumptions on the weights \cite{D04,DG21,DSWW06,K03,NC06,SR02}. Randomized lattice rules further yield sharp root-mean-square error bounds and often improve robustness with respect to unknown problem structure \cite{DGS22,G26,GK26,KKNU19,KNW23}.

Beyond integration, lattice sampling has become an increasingly important paradigm for \emph{approximation} and \emph{reconstruction} of multivariate periodic functions. Early investigations demonstrated how lattice information can be used to recover trigonometric polynomials and to design efficient spectral algorithms \cite{KSW06,LH03}, see also \cite{HW63}. More recent work established systematic lattice-based approximation schemes in periodic spaces with general weights together with fast CBC constructions for approximation settings \cite{CKNS20,CKNS21,KMNN21}. Complementary strands of research analyze approximation based on single or multiple rank-1 lattices from a harmonic analysis viewpoint, emphasizing sampling stability and explicit discretizations for multivariate trigonometric polynomials \cite{BKTV17,K19,KPV15,KV19}. The continuing interest in this topic is also reflected in recent advances on subsampled and structured lattices aimed at optimal $L_2$ approximation error convergence \cite{BGKS24}.

A key challenge in lattice-based approximation is to control aliasing and to obtain \emph{high-probability} guarantees for randomized constructions with minimal overhead. In the wider randomized numerical analysis and QMC literature, the \emph{median} and \emph{median-of-means} principles have emerged as effective tools for turning average-case or second-moment bounds into sharp tail estimates without requiring precise knowledge of smoothness parameters. This idea goes back to robust estimation in Monte Carlo contexts \cite{KNR19,NP09} and has recently been used to construct smoothness-adaptive or nearly universal randomized QMC strategies for integration \cite{GK26,GL22,GSM24,P26,PO23,PO24}. These developments suggest that median aggregation is a natural candidate for stabilizing lattice-based \emph{approximation} algorithms as well.

The present note fits into this emerging line of research. Building on the $L_2$-approximation analysis of median lattice algorithms in Korobov spaces from \cite{PKG25} (see also \cite{PGK25}) and the recent multiple-shift lattice approximation framework of \cite{CDG25}, we investigate the worst-case $L_p$-approximation for the full range $1 \le p \le \infty$. We consider a randomized median lattice algorithm that approximates Fourier coefficients on a weighted hyperbolic-cross-type index set by $R$ independent rank-1 lattice rules with a prime number of points, and then applies the componentwise median across the repetitions. The main result shows that this approach yields, with high probability, nearly optimal convergence rates in $L_p$, with the expected logarithmic factor associated with hyperbolic-cross truncation and an arbitrarily small loss parameter $\beta>0$. In the endpoint case $p=\infty$, we additionally obtain dimension-independent constants under standard summability conditions on the product weights, mirroring the behavior known from related integration and approximation settings in weighted Korobov spaces \cite{DKP22,KKNU19,KMNN21}.



Our main result is as follows.
\begin{theorem}\label{thm:main}
Let $\alpha>1/2$ and $d \in \NN$, and let $\gamma_j \in (0,1]$ for $j \in \{1, \ldots, d\}$.
Let $A$ be the approximation algorithm given by Algorithm~\ref{alg:median} with a parameter $\tau>0$, an odd number $R>1 $, and a prime number $N$. Let $\varepsilon_2$ be given by \eqref{eqn:failprob2} and assume it is at most 1.
Then for any $1 \le p \le \infty$, with probability at least $1 - \varepsilon_2$, the worst-case $L_p$-approximation error of $A$ for function approximation in the weighted periodic Korobov space $\Hcal_{d, \alpha, \bsgamma}$ given by \eqref{Korspace} satisfies
$$    \mathrm{err}(\Hcal_{d,\alpha,\bsgamma}, L_p, A) \leq     C_{d, \alpha, \beta, \bsgamma, \tau, R,p}\,  N^{ - \alpha + (\frac12 - \frac1p)_+ + \beta },$$ for any $\beta > 0$, where $(x)_+ = \max\{x, 0\}$, $C_{d, \alpha, \beta, \bsgamma, R,\tau,p} > 0$ is a constant independent of $N$ but dependent on $d$. For $p= \infty$, the constant is independent of the dimension if $\sum_{j=1}^\infty \gamma_j^{1/(2\alpha)} < \infty$.
\end{theorem}

\begin{remark}
Theorem \ref{thm:main} is proven in Section~\ref{sec_Lp}. In particular, it  matches the main, optimal convergence rates for linear sampling algorithms, see \cite[Section~5.1 and Corollary~19]{KPUU25}. We note that it follows from \cite[Theorem~4.7]{NW08} that non-adaptive linear algorithms are optimal for the problems considered here.
\end{remark}

\medskip

The paper is organized as follows. Section~\ref{sec:preliminaries} introduces notation, lattice rules, weighted Korobov spaces, and the approximation error criterion. In Section~\ref{sec:median}, we describe the median lattice algorithm. Section~\ref{sec:main} contains the high-probability error analysis in $L_\infty$ and $L_2$. The proof of Theorem~\ref{thm:main} is based on the $L_2$-approximation and $L_\infty$-approximation error bounds together with a norm interpolation inequality, and is presented in Section~\ref{sec_Lp}.





\section{Preliminaries and notation}\label{sec:preliminaries}


Throughout, let \(d\in \NN\), where $\NN$ is the set of positive integers, be the dimension and let boldface symbols denote vectors in \(\RR^d\) or \(\ZZ^d\). 


For a multi-index \(\bsh=(h_1,\ldots,h_d)\in\ZZ^d\), its support is denoted by
\[
\supp(\bsh):=\{\,j\in\{1,\ldots,d\}\,:\, h_j\neq 0\,\}.
\]

Let \(\bsgamma=(\gamma_{\mathfrak u})_{\mathfrak u\subseteq\{1,\ldots,d\}}\) be a family of positive weights with \(\gamma_{\emptyset}=1\). For a smoothness parameter \(\alpha>0\), define the frequency weight
\[
r_{2\alpha,\bsgamma} (\bsh):=\gamma_{\supp(\bsh)}^{-1}\prod_{j\in \supp(\bsh)} \abs{h_j}^{2\alpha},
\qquad \bsh\in\ZZ^d,
\]
so that in particular \(r_{2\alpha,\bsgamma}(\boldsymbol 0)=1\).


For \(f\in L_2([0,1]^d)\) we use the Fourier coefficients
\[
\widehat{f}(\bsh)\;:=\;\int_{[0,1]^d} f(\bsx)\,e^{-2\pi i\,\bsh\cdot\bsx}\, \mathrm d\bsx,
\qquad \bsh\in\ZZ^d,
\]
where \(\bsh\cdot\bsx=\sum_{j=1}^d h_j x_j\).

The weighted periodic Korobov space \(\Hcal_{d,\alpha,\bsgamma}\) is the collection of all \(f\in L_2([0,1]^d)\) with finite norm
\[
  \|f\|_{d,\alpha,\bsgamma}^2:=\sum_{\bsh\in\ZZ^d} \abs{\widehat{f}({\bsh})}^2 r_{2\alpha,\bsgamma}(\bsh),
\]
i.e.,
\begin{equation}\label{Korspace}
\Hcal_{d,\alpha, \bsgamma} = \{ f: [0,1]^d \to \RR\,:\, \|f\|_{d, \alpha, \bsgamma} < \infty\},
\end{equation}
where the norm arises from the inner product
\[
\langle f,g\rangle_{d,\alpha,\bsgamma}\;:=\;\sum_{\bsh\in\ZZ^d} \overline{\widehat f(\bsh)}\,\widehat g(\bsh)\, r_{2\alpha,\bsgamma}(\bsh),
\]
under which $\Hcal_{d,\alpha,\bsgamma}$ is a Hilbert space.


For \(g: [0,1]^d\to \RR \) we use the standard Lebesgue norms
\[
\|g\|_{L_p}:=\left(\int_{[0,1]^d}\abs{g(\bsx)}^{p}\,\mathrm d\bsx\right)^{1/p},
\qquad
\|g\|_{L_\infty}:=\operatorname*{ess\,sup}_{\bsx\in [0,1]^d}\abs{g(\bsx)}.
\]



An approximation algorithm  $A$ is a map $A:\Hcal_{d,\alpha,\bsgamma}\to L_p( [0,1]^d)$ (for $1 \le p \le \infty $) that assigns to each $f$ an approximation 
$A(f)$. The worst-case error of $A$ over the unit ball of 
$\Hcal_{d,\alpha,\bsgamma}$ is measured by
\[
  \mathrm{err}(\Hcal_{d,\alpha,\bsgamma}, L_p, A):= \sup_{\substack{f\in \Hcal_{d,\alpha,\bsgamma}\\ 
  \norm{f}_{d,\alpha,\bsgamma}\le 1}}
  \|f-A(f)\|_{L_p}, \quad 1 \le p \le \infty.
\]

In the following, we will frequently write $\{k{:}\ell\}$ to denote the set $\{k,k+1,k+2,\ldots,\ell\}$ for integers $k\le \ell$.
The role of $\bsgamma$ is to modulate the importance of coordinate subsets (see \cite{SW98}): frequencies supported on $\mathfrak u\subseteq\{1{:}d\}$ are scaled by $\gamma_{\mathfrak u}^{-1}$. In the present paper, we consider the most prominent class, namely product weights $\gamma_{\mathfrak u}=\prod_{j\in\mathfrak u}\gamma_j$ with $\gamma_{\emptyset}=1$, where $(\gamma_j)_{j\ge 1}$ is a non-increasing sequence of positive reals. We remark that we could also allow non-negative $\gamma_j$, but this would result in additional technical notation, which we avoid. The factor $\prod_{j\in\supp(\bsh)} |h_j|^{2\alpha}$ in the definition of $r_{2\alpha,\bsgamma}$, and thus also in the norm of $\Hcal_{d,\alpha,\bsgamma}$, penalizes higher frequencies according to the smoothness parameter $\alpha$; larger $\alpha$ corresponds to stronger decay of Fourier coefficients in $\Hcal_{d,\alpha,\bsgamma}$. 

Let us now outline the definition of the QMC rules we use in our median algorithm, namely rank-1 lattice rules. For simplicity, let us assume that $N$ is a prime in the following (assuming that $N$ is prime guarantees that all one-dimensional projections of a lattice point set are evenly distributed, which is an advantage. If we would allow composite $N$, we presumably could derive similar results, but would need more technical notation). Let $\bsz=(z_1,\ldots,z_d)$ be a 
vector with each component in $\{1{:}(N-1)\}$. We can then define a 
lattice point set $P_{N,\bsz}$ with points $\bsx_0,\ldots,\bsx_{N-1}$ as follows. For $j\in\{1{:}d\}$ and $k\in\{0,\ldots,N-1\}$, the $j$-th component of $\bsx_k$ is given by
\[
x_k^{(j)}:=\left\{\frac{k z_j}{N}\right\},
\]
where $\{y\}=y-\lfloor y \rfloor$ denotes the fractional part of a real number $y$. By using this definition for all $j\in\{1{:}d\}$ and $k\in\{0,\ldots,N-1\}$, we obtain the full lattice point set $P_{N,\bsz}$. A QMC rule using $P_{N,\bsz}$ is called a (rank-1) lattice rule. We further define the dual lattice of $P_{N,\bsz}$ by
\begin{align}\label{eq:def_dual_lattice}
P_{N,\bsz}^\perp := \{\bsl \in \ZZ^d \,:\,  \bsz^\top \bsl \equiv 0 \tmod N\}.
\end{align}

Note that, for fixed $N$ and $d$, a lattice rule is fully characterized by the choice of the \textit{generating vector} $\bsz$. Not all choices of $\bsz$ yield lattice rules of sufficient quality as integration node sets. However, there are fast construction algorithms available that return good generating vectors for given $N$, $d$, $\alpha$, and $\bsgamma$. We refer to the books \cite{DKP22, SJ94} for overviews of the theory of lattice rules, and in particular to \cite[Chapters 3 and 4]{DKP22} for constructions of good lattice rules. In the present paper, however, we will make a random choice of the generating vectors $\bsz$ to obtain the lattice point sets used. Therefore, we need not be concerned with the search for good generating vectors here, which is a general computational advantage of median rules. 

It is sometimes useful to introduce an additional random element when applying lattice rules. This is usually done by a \emph{random shift}, $\Delta\in [0,1)^d$. Given $P_{N,\bsz}=\{\bsx_0,\ldots,\bsx_{N-1}\}$ and drawing $\bsDelta$ from a uniform distribution over $[0,1)^d$, the corresponding \emph{randomly shifted lattice rule} applied to a function $g$ is then 
\[
   Q_{N,d,\bsDelta}(g):=\frac{1}{N}\sum_{k=0}^{N-1}g \left(\left\{\bsx_k + \bsDelta\right\}\right),
\]
i.e., all points of $P_{N,\bsz}$ are shifted modulo one by the same $\bsDelta$.


\section{The median algorithm}\label{sec:median}

For a real $L\geq 0$, we put 
\begin{equation}\label{eq:Acal}
\Acal_{d,\alpha,\bsgamma}(L) := \left\{ \bsh \in \ZZ^d \,:\,  r_{2\alpha, \bsgamma}(\bsh) < L^{2\alpha} \right\}=\left\{ \bsh \in \ZZ^d \,:\,  r_{1, \bsgamma^{1/(2\alpha)}}(\bsh) < L\right\}.
\end{equation}

For an odd positive integer $R$, we define the median of $R$ complex numbers $Z_1,\dots,Z_R$ as
\begin{equation}\label{eqn:complexmedian}
    \mathrm{median}_{r\in\{1{:}R\}}(Z_r)=\mathrm{median}_{r\in\{1{:}R\}} \Re(Z_r) + \icomp\cdot \mathrm{median}_{r\in\{1{:}R\}} \Im(Z_r),
\end{equation}
where $\Re(Z_r)$ and $\Im(Z_r)$ denote the real and imaginary parts of $Z_r$, respectively.
In the following, we will also assume that $N$ is an odd prime 
number. Though we do not need this requirement for all results below, we sometimes do, and it is therefore easier to make this assumption from now on.
\begin{algorithm}\label{alg:median}
    Let $\alpha>1/2$ and $d \in \NN$, and let $\gamma_j \in (0,1]$ for 
    $j \in \{1, \ldots, d\}$. Given  a parameter $\tau > 0$, an odd number of repetitions $R >1$, and a prime number $N$, do the following.
    \begin{enumerate}
        \item Define 
        \begin{equation}\label{eqn:PN}
        P_{N,d}(\tau,\alpha,\bsgamma):=\prod_{j=1}^d \left(1+ 2\gamma_j^{1/(2\alpha)} (1+\tau\log N)\right),
        \end{equation}
        and
        \begin{equation}\label{eqn:Nstar}
        N_2:=\frac{N-1 }{\exp(\tau^{-1})P_{N,d}(\tau,\alpha,\bsgamma)}.
        \end{equation}
        \item \textbf{For} $r$ from $1$ to $R$, do the following:
        \begin{enumerate}
            \item Randomly draw $\bsz_{r}$ from the uniform distribution over the set $\{1{:}(N-1)\}^d$.
            \item For all $\bsh\in \Acal_{d,\alpha,\bsgamma}(N_2)$, set
            \[ 
            \widehat{f}_{N,\bsz_r}(\bsh):=
            \frac{1}{N} \sum_{k=0}^{N-1} f\left( \left\{ \frac{k \bsz_r}{N} \right\}\right)\exp\left(-2\pi \icomp \bsh \cdot \left( \frac{k \bsz_r}{N}\right)\right).
            \]   
        \end{enumerate}
       \textbf{end for}
       \item For all $\bsh\in \Acal_{d,\alpha,\bsgamma}(N_2)$, define the aggregation
        \[ 
            \widehat{f}_{N,\bsz_{\{1{:}R\}}}(\bsh):=\mathrm{median}_{r\in\{1{:}R\}} \Big(\widehat{f}_{N,\bsz_r}(\bsh)\Big),
       \]
       where $\bsz_{\{1{:}R\}}$ denotes the collections of $\bsz_r$ for $r\in\{1{:}R\}$, respectively.
       Define the final approximation by
        \[ 
        (A_{N,\bsz_{\{1{:}R\}},\Acal_{d,\alpha,\bsgamma}(N_2)}(f))(\bsx):= \sum_{\bsh \in \Acal_{d,\alpha,\bsgamma}(N_2)} \widehat{f}_{N,\bsz_{\{1{:}R\}}}(\bsh) \exp (2 \pi \icomp \bsh \cdot \bsx).\]
\end{enumerate}
\end{algorithm}

\begin{remark}
    The parameter $\tau>0$ does not affect the asymptotic convergence rates of Algorithm~\ref{alg:median} to be established in Section~\ref{sec:main}. For guidance on optimal tuning of $\tau$ to improve finite-sample performance, see \cite[Remark 16]{PKG25}.
\end{remark}

\begin{remark}
%Our previous median lattice algorithm \cite{PKG25} differs slightly from Algorithm~\ref{alg:median} in that it 
A natural modification of Algorithm~\ref{alg:median} is to incorporate additional independent uniform shifts $\bsDelta_1,\dots,\bsDelta_R \in [0,1)^d$, and replace $\widehat{f}_{N,\bsz_{\{1{:}R\}}}(\bsh)$ by 
       \[ 
            \widehat{f}_{N,\bsz_{\{1{:}R\}},\bsDelta_{\{1{:}R\}}}(\bsh):=\mathrm{median}_{r\in\{1{:}R\}} \Big(\widehat{f}_{N,\bsz_r,\bsDelta_r}(\bsh)\Big),
       \]
       for 
                   \[ 
            \widehat{f}_{N,\bsz_r,\bsDelta_r}(\bsh):=
            \frac{1}{N} \sum_{k=0}^{N-1} f\left( \left\{ \frac{k \bsz_r}{N}+\bsDelta_r \right\}\right)\exp\left(-2\pi \icomp \bsh \cdot \left( \frac{k \bsz_r}{N}+\bsDelta_r\right)\right).
            \]  
            We will study Algorithm~\ref{alg:median} for simplicity, but all error bounds stated in Section~\ref{sec:main} also hold under the above change due to the common upper bounds
            $$\abs{\widehat{f}_{N,\bsz_r}(\bsh)-\widehat{f}({\bsh})}=\abs{\sum_{\bsl\in P^\perp_{N,\bsz_r}\setminus\{\bszero\} }\widehat{f}({\bsh+\bsl})}\leq \sum_{\bsl\in P^\perp_{N,\bsz_r}\setminus\{\bszero\} }\abs{\widehat{f}({\bsh+\bsl})} $$
            and 
            $$\abs{\widehat{f}_{N,\bsz_r,\bsDelta_r}(\bsh)-\widehat{f}({\bsh})}=\abs{\sum_{\bsl\in P^\perp_{N,\bsz_r}\setminus\{\bszero\} }\widehat{f}({\bsh+\bsl})\exp(2\pi\icomp \bsl\cdot\bsDelta_r)}\leq \sum_{\bsl\in P^\perp_{N,\bsz_r}\setminus\{\bszero\} }\abs{\widehat{f}({\bsh+\bsl})},$$
         where the dual lattice $P^\perp_{N,\bsz_r}$ is given by \eqref{eq:def_dual_lattice}.
            
    When measured in the root-mean-squared error, which is the setting considered in \cite{PKG25}, $\widehat{f}_{N,\bsz_r,\bsDelta_r}(\bsh)$ enjoys a tighter error bound
    \begin{equation*}
      \left(\int_{[0,1]^d} |\widehat{f}_{N,\bsz_r,\bsDelta_r}(\bsh)-\widehat{f}(\bsh)|^2 \rd\bsDelta_r\right)^{1/2}= \left(\sum_{\bsl \in P_{N,\boldsymbol{z}_r}^\perp\setminus \{\bszero\}}|\widehat{f}(\bsh+\bsl)|^2\right)^{1/2}.
    \end{equation*}
    Therefore, we expect median lattice algorithms with the random shifts to perform better than those without the shifts, despite the two sharing the same worst-case error bound.
\end{remark}



\section{Approximation error bounds}\label{sec:main}

Throughout this section, we assume that $\tau>0$, $R$ is an odd number greater than $1$, and $N$ is a sufficiently large prime number such that $N_2> 1$.

\subsection{$L_\infty$-approximation}

The next lemma provides a bound on the size of the frequency set $\Acal_{d,\alpha, \bsgamma}(N_2)$, which was shown in \cite[Corollary 6]{PKG25}.
\begin{lemma}\label{lem:Adcount} For $N_2$ defined as in \eqref{eqn:Nstar} and $\Acal_{d,\alpha,\bsgamma}(N_2)$ as in \eqref{eq:Acal}, 
we have
\[ 
 |\Acal_{d,\alpha,\bsgamma}(N_2)|\leq 1+\frac{ N-1}{1+\tau\log N_2}.
\]
\end{lemma}

For a given generating vector $\bsz$, we define the set $K_{d, \alpha, \bsgamma}(\bsz)$ of frequencies $\bsh \in \Acal_{d, \alpha, \bsgamma}$ for which no frequency $\bsk \in \Acal_{d, \alpha, \bsgamma}$ is aliased with $\bsh$, i.e., we have $(\bsh - \bsk) \cdot \bsz \not\equiv 0 \pmod{N}$ for all $\bsk \in \Acal_{d, \alpha, \bsgamma}\setminus \{\bsh\}$: 
\begin{equation*}
K_{d,\alpha,\bsgamma}(\bsz):=\left\{\bsh\in \Acal_{d,\alpha,\bsgamma}(N_2)\,:\,   (\bsh+ P^\perp_{N,\bsz}) \bigcap \Acal_{d,\alpha,\bsgamma}(N_2)=\{\bsh\} \right\}.    
\end{equation*}
This set is crucial in the proof of our main result, since for ``good'' generating vectors $\bsz$ we obtain ``good'' approximations of the Fourier coefficients for frequencies in $K_{d,\alpha, \bsgamma}(\bsz)$. Roughly speaking, the proof strategy is to show that for each frequency $\bsh \in \Acal_{d,\alpha, \bsgamma}(N_2)$, more than half of the $\bsz_1, \ldots, \bsz_R$ are such that $\bsh \in K_{d,\alpha, \bsgamma}(\bsz_r)$.

Next, we show several auxiliary results.
\begin{lemma}\label{lem:reconstruction}
    For $\bsh\in \Acal_{d,\alpha,\bsgamma}(N_2) $ and $\bsz$ chosen uniformly from  $\{1{:}(N-1)\}^d$, 
    $$\Pr(\bsh\notin K_{d,\alpha,\bsgamma}(\bsz) )\leq \frac{1}{1+\tau\log N_2}.$$
\end{lemma}
\begin{proof}
    By Lemma~\ref{lem:Adcount}, we have $|\Acal_{d,\alpha,\bsgamma}(N_2)|<N$. Consequently, the one-dimensional projection of $\Acal_{d,\alpha,\bsgamma}(N_2)$ onto any $j$-th coordinate, defined as
    \[ \proj_j (\Acal_{d,\alpha,\bsgamma}(N_2)) := \left\{ h_j\in \ZZ \,:\,  \bsh\in \Acal_{d,\alpha,\bsgamma}(N_2) \right\}, \]
    is contained in the set $\{-(N-1)/2,\ldots,0,\ldots,(N-1)/2\}$. Thus, for any distinct elements $\bsh,\bsh'\in \Acal_{d,\alpha,\bsgamma}(N_2)$, it follows that $h_j-h'_j \not\equiv 0 \pmod N$ for at least one $j$. This implies that $\bsh\not\equiv \bsh' \pmod N$, or equivalently, $N\nmid \bsh-\bsh'$.
    
    
    Now, for each $\bsh'\in \Acal_{d,\alpha,\bsgamma}(N_2) $ distinct from $\bsh$, the fact that $N\nmid \bsh-\bsh'$ leads to
    $$\Pr\left(\bsz\cdot (\bsh'-\bsh)\equiv 0 \tmod N  \right)\leq \frac{1}{N-1}.$$
    Hence, $\Pr(\bsh'-\bsh\in P^\perp_{N,\bsz})\leq 1/(N-1)$. By a union bound, we obtain
    $$\Pr(\bsh\notin K_{d,\alpha,\bsgamma}(\bsz))\leq \frac{| \Acal_{d,\alpha,\bsgamma}(N_2)|-1}{N-1}\leq\frac{1}{1+\tau\log N_2}.$$
\end{proof}

\begin{corollary}\label{cor:whp}
    Let 
    \begin{equation}\label{eqn:failprob}
      \varepsilon_1:=\frac{|\Acal_{d,\alpha,\bsgamma}(N_2)|}{2}\left(\frac{4}{1+\tau\log N_2}\right)^{\lceil R/2 \rceil}.  
    \end{equation}
    If $ \varepsilon_1\leq 1$, then with probability at least $1- \varepsilon_1$, 
    $$
    \min_{\bsh\in \Acal_{d,\alpha,\bsgamma}(N_2)} \sum_{r=1}^R \bsone\{\bsh\in K_{d,\alpha,\bsgamma}(\bsz_r)\}\geq \lceil R/2 \rceil,
    $$
    where $\bsone$ denotes the indicator function of a given set.
\end{corollary}
\begin{proof}
    For each $\bsh\in \Acal_{d,\alpha,\bsgamma}(N_2) $, the probability that $\bsh$ does not belong to at least $\lceil R/2 \rceil$ instances of $K_{d,\alpha,\bsgamma}(\bsz_r)$ among $\{K_{d,\alpha,\bsgamma}(\bsz_1),\dots,K_{d,\alpha,\bsgamma}(\bsz_R)\}$ is bounded by
    $$2^R\left(\frac{1}{1+\tau\log N_2}\right)^{\lceil R/2 \rceil}=\frac{1}{2} \left(\frac{4}{1+\tau\log N_2}\right)^{\lceil R/2 \rceil} .$$
    The conclusion follows from a union bound over $\bsh\in \Acal_{d,\alpha,\bsgamma}(N_2) $.
\end{proof}

\begin{lemma}\label{lem:medtosum}
    Let $\{x_1,\dots,x_R\}\subseteq [0,\infty)$. For any $J\subseteq \{1{:}R\}$ with $|J|\geq \lceil R/2 \rceil$,
    $$\mathrm{median}_{r\in\{1{:}R\}} x_r \leq \sum_{r\in J} x_r.$$
\end{lemma}
\begin{proof}
    Let $r^*\in \{1{:}R\}$ be such that it satisfies $x_{r*}=\mathrm{median}_{r\in\{1{:}R\}} x_r$. If $r^*\in J$, the inequality is trivial. When $r^*\notin J$, we must have
    $$x_{r*}\leq \max_{r\in J} x_r\leq \sum_{r\in J} x_r. $$
\end{proof}

We need the following lemma, whose proof can be found, for instance, in \cite[Chapter~14, Lemma~14.2]{DKP22}.
\begin{lemma}\label{lem_sum}
Let $\alpha>1/2$ and $M\geq 1$ be given. For any $q\in(1/(2\alpha),1)$ we have 
\begin{align*}
\sum_{\bsk\in\mathbb{Z}^d\backslash \Acal_{d, \alpha,\bsgamma}(M)}\frac{1}{r_{2\alpha,\bsgamma}(\bsk)}\leq \frac{1}{(\gamma_1^{1/2} M^{2\alpha})^{(1/q-1)/(2\alpha)}}\frac{q}{1-q}\prod_{j=1}^d\left(1+2\gamma_j^{q} \zeta(2\alpha q)  \right)^{1/q}.
\end{align*}
\end{lemma}

\begin{theorem}\label{thm:infty}
Let $\varepsilon_1$ be given by \eqref{eqn:failprob} and assume it 
is at most 1. Let $\alpha>1/2$ and $M\ge 1$ be given.
Then, for any $q\in(1/(2\alpha),1)$, with probability at least $1-\varepsilon_1$,
    $$\mathrm{err}(\Hcal_{d,\alpha,\bsgamma}, L_\infty, A_{N,\bsz_{\{1{:}R\}},\Acal_{d,\alpha,\bsgamma}(N_2)})\leq ( 2R+1) \left( \frac{1}{(\gamma_1^{1/(4\alpha)} N_2)^{(1/q-1)}}\frac{q}{1-q}\prod_{j=1}^d\left(1+2\gamma_j^{q} \zeta(2\alpha q)  \right)^{1/q} \right)^{1/2}. % \left(\sum_{\bsh\in \ZZ^d\setminus\Acal_{d,\alpha,\bsgamma}(N_2)} \frac{1}{r_{2\alpha,\bsgamma}(\bsh)}\right)^{1/2}. 
    $$
\end{theorem}

\begin{proof}
    For $f\in \Hcal_{d,\alpha,\bsgamma}$,
    $$\Vert A_{N,\bsz_{\{1{:}R\}},\Acal_{d,\alpha,\bsgamma}(N_2)}(f)-f\Vert_{L_\infty}\leq \sum_{\bsh\in \Acal_{d,\alpha,\bsgamma}(N_2)} \abs{\widehat{f}_{N,\bsz_{\{1{:}R\}}}(\bsh)-\widehat{f}({\bsh})}+\sum_{\bsh\in \ZZ^d\setminus\Acal_{d,\alpha,\bsgamma}(N_2)} \abs{\widehat{f}({\bsh})}.$$

In the following we use that for a set of real numbers $x_1, x_2, \ldots, x_R$, with $R$ odd, we have $| \mathrm{median}_{r \in \{1:R\}} x_r | \le \mathrm{median}_{r \in \{1:R\}} |x_r|$. Then the complex median \eqref{eqn:complexmedian} satisfies
\begin{align*}
    \abs{\widehat{f}_{N,\bsz_{\{1{:}R\}}}(\bsh)-\widehat{f}({\bsh})} & \leq \mathrm{median}_{r\in\{1{:}R\}} \abs{\Re\left(\widehat{f}_{N,\bsz_{r}}(\bsh)-\widehat{f}({\bsh})\right)}+\mathrm{median}_{r\in\{1{:}R\}} \abs{\Im\left(\widehat{f}_{N,\bsz_{r}}(\bsh)-\widehat{f}({\bsh})\right)}\\
    & \leq 2 \ \mathrm{median}_{r\in\{1{:}R\}} \abs{\widehat{f}_{N,\bsz_{r}}(\bsh)-\widehat{f}({\bsh})}.
\end{align*}
By Corollary~\ref{cor:whp} and Lemma~\ref{lem:medtosum}, with probability at least $1-\varepsilon_1$, it holds that
    \begin{align*}
        \sum_{\bsh\in \Acal_{d,\alpha,\bsgamma}(N_2)} \abs{\widehat{f}_{N,\bsz_{\{1{:}R\}}}(\bsh)-\widehat{f}({\bsh})} & \leq 2\sum_{\bsh\in \Acal_{d,\alpha,\bsgamma}(N_2)}  \mathrm{median}_{r\in\{1{:}R\}} \abs{\widehat{f}_{N,\bsz_r}(\bsh)-\widehat{f}({\bsh})} \\
        & \leq 2 \sum_{r=1}^R \sum_{\bsh\in K_{d,\alpha,\bsgamma}(\bsz_r)} \abs{\widehat{f}_{N,\bsz_r}(\bsh)-\widehat{f}({\bsh})}\\
        & \leq 2 \sum_{r=1}^R \sum_{\bsh\in K_{d,\alpha,\bsgamma}(\bsz_r)} \sum_{\bsl\in P^\perp_{N,\bsz_r}\setminus\{\bszero\} }\abs{\widehat{f}({\bsh+\bsl})}\\
        & \leq 2R \sum_{\bsh\in \ZZ^d\setminus\Acal_{d,\alpha,\bsgamma}(N_2)} \abs{\widehat{f}({\bsh})},
    \end{align*}
    where the last inequality follows because $(\bsh,\bsl) \to \bsh+\bsl$ is injective over $K_{d,\alpha,\bsgamma}(\bsz_r)\times  P^\perp_{N,\bsz_r} $, and $\bsh+\bsl\in \Acal_{d,\alpha,\bsgamma}(N_2)$ holds only if $\bsl=\bszero$. Therefore, using the Cauchy-Schwarz inequality,
    \begin{align*}
        \Vert A_{N,\bsz_{\{1{:}R\}},\Acal_{d,\alpha,\bsgamma}(N_2)}(f)-f\Vert_{L_\infty} & \leq  (2 R+1)\sum_{\bsh\in \ZZ^d\setminus\Acal_{d,\alpha,\bsgamma}(N_2)} \abs{\widehat{f}({\bsh})} \\
        & \leq (2 R+1) \|f\|_{d,\alpha,\bsgamma} \left(\sum_{\bsh\in \ZZ^d\setminus\Acal_{d,\alpha,\bsgamma}(N_2)} \frac{1}{r_{2\alpha,\bsgamma}(\bsh)}\right)^{1/2}.
    \end{align*}
    The conclusion follows immediately from Lemma~\ref{lem_sum}.
\end{proof}
Since $q$ can be arbitrarily close to $1/(2\alpha)$, $\mathrm{err}(\Hcal_{d,\alpha,\bsgamma}, L_\infty, A_{N,\bsz_{\{1{:}R\}},\Acal_{d,\alpha,\bsgamma}(N_2)})$ is nearly $O(N^{-\alpha+1/2})$ with probability $1-\varepsilon_1$. The bound is independent of the weights if $\sum_{j=1}^\infty \gamma_j^{1/(2\alpha)} < \infty$.



\subsection{$L_2$-approximation}

For $\bsl\in \ZZ^d$, we let
$$\rho_{2\alpha, \bsgamma}(\bsz)=\min_{\bsl\in P^\perp_{N,\bsz}\setminus\{\bszero\}} r_{2\alpha,\bsgamma}(\bsl).$$


The following result is shown in \cite[Eq. (3.9)]{CDG25}. We state it as a lemma.
\begin{lemma}
Let $\bsz \in \{1{:}(N-1)\}^d$. Let $\bsgamma = (\gamma_1, \gamma_2, \ldots)$ with $\gamma_1 \ge \gamma_2 \ge \cdots$. Let $Q > 1$, $\beta>1/2$, and $\rho_{\beta, \bsgamma}(\bsz) > 1$. Let $m \in \mathbb{Z}$ be such that $2^{m-1} < M\le 2^m$. Then
\begin{equation*}
\sup_{\bsh \in \mathbb{Z}^d} \left| \bsl \in \bsh + P^\perp_{N, \bsz} : r_{\beta, \bsgamma}(\bsl) \le Q \right|   \le \frac{2^m}{m} \min \left\{ \frac{(1 + m/2)^d}{\rho_{\beta, \bsgamma}(\bm{z})^{1/\beta}}, \prod_{j=1}^d (1 + \gamma_j^{1/\beta} m)  \right\}.
\end{equation*}
\end{lemma}
In the notation of the present paper, where $\beta$ is replaced by $2 \alpha$, $Q = L^{2\alpha}$, and hence $2^{m-1} < L \le 2^m$, this reads
\begin{equation}\label{fiber:bound}
\sup_{\bsh \in \mathbb{Z}^d} \left| \bsl \in \bsh + P^\perp_{N, \bsz} : r_{2 \alpha, \bsgamma}(\bsl) \le L^{2\alpha}  \right|  \le \frac{2^m}{m} \min \left\{ \frac{(1 + m/2)^d}{\rho_{2 \alpha, \bsgamma}(\bm{z})^{1/(2\alpha) }}, \prod_{j=1}^d (1 + \gamma_j^{1/(2\alpha) } m)  \right\}.
\end{equation}

For $L_\infty$-approximation it was sufficient to show that for each $\bsh \in \Acal_{d,\alpha, \bsgamma}(N_2)$, more than half of the generating vectors satisfy $\bsh \in K_{d,\alpha, \bsgamma}(\bsz_r)$, which was shown in Corollary~\ref{cor:whp}. However, for $L_2$-approximation we need the additional property that for all generating vectors the figure of merit satisfies $\rho_{2\alpha, \bsgamma}(\bsz_r) \ge N_2^{2\alpha}$, which we show in the following lemma. Ensuring that the figure of merit is sufficiently large will allow us to improve the convergence rate of the $L_2$-approximation error by a factor of $N_2^{1/2}$ compared to the $L_\infty$-approximation error.

\begin{lemma}\label{cor:whp2}
    Let 
    \begin{equation}\label{eqn:failprob2}
      \varepsilon_2:=\frac{|\Acal_{d,\alpha,\bsgamma}(N_2)|}{2}\left(\frac{8}{1+\tau\log N_2}\right)^{\lceil R/2 \rceil}.  
    \end{equation}
    If $\varepsilon_2\leq 1$, then with probability at least $1- \varepsilon_2$, 
    \[
      \min_{\bsh\in \Acal_{d,\alpha,\bsgamma}(N_2)} \sum_{r\in \mathcal{R}} \bsone\{\bsh\in K_{d,\alpha,\bsgamma}(\bsz_r)\}\geq \lceil R/2 \rceil
    \]
    for 
    \[
      \mathcal{R}=\left\{r\in \{1{:}R\} \,:\, \rho_{2\alpha, \bsgamma}(\bsz_r)\geq N^{2\alpha}_2 \right\}.
    \]
\end{lemma}
\begin{proof}
For any $r\in \{1{:}R\}$, when $\bszero\in K_{d,\alpha,\bsgamma}(\bsz_r)$, no $\bsl\in P^\perp_{N,\bsz_r}$ other than $\bszero$ belong to $\Acal_{d,\alpha,\bsgamma}(N_2)$, so
$$\rho_{2\alpha, \bsgamma}(\bsz_r)=\min_{\bsl\in P^\perp_{N,\bsz_r}\setminus\{\bszero\}} r_{2\alpha,\bsgamma}(\bsl)\geq N^{2\alpha}_2$$
and $r\in\mathcal{R}$. Consequently, for any $\bsh\in  \Acal_{d,\alpha,\bsgamma}(N_2)$,
\[
  \sum_{r\in \mathcal{R}} \bsone\{\bsh\in K_{d,\alpha,\bsgamma}(\bsz_r)\} \ge 
  \sum_{r\in \{1:R\}} \bsone\{\bszero\in K_{d,\alpha,\bsgamma}(\bsz_r)\text{ and }\bsh\in K_{d,\alpha,\bsgamma}(\bsz_r)\}.
\]
The probability that
\[
\sum_{r\in \{1:R\}} \bsone\{\bszero\in K_{d,\alpha,\bsgamma}(\bsz_r)\text{ and }\bsh\in K_{d,\alpha,\bsgamma}(\bsz_r)\} < \lceil R/2 \rceil 
\]
for at least one $\bsh\in \Acal_{d,\alpha,\bsgamma}(N_2)$ is equivalently the probability that 
\begin{equation}\label{eq:probability_converse}
 \sum_{r\in \{1:R\}} \bsone\{\bszero\notin K_{d,\alpha,\bsgamma}(\bsz_r)\text{ or }\bsh\notin K_{d,\alpha,\bsgamma}(\bsz_r)\} \ge \lceil R/2 \rceil .
\end{equation}
for at least one $\bsh\in \Acal_{d,\alpha,\bsgamma}(N_2)$.
However, Lemma~\ref{lem:reconstruction} implies for every $r\in \{1{:}R\}$ and $\bsh\in  \Acal_{d,\alpha,\bsgamma}(N_2)$,
$$\Pr\left(\bszero\notin K_{d,\alpha,\bsgamma}(\bsz_r)\text{ or } \bsh\notin K_{d,\alpha,\bsgamma}(\bsz_r)\right)\leq \frac{2}{1+\tau\log N_{2}}.$$
Therefore, the probability that \eqref{eq:probability_converse} holds for a given $\bsh\in \Acal_{d,\alpha,\bsgamma}(N_2)$ is bounded from above by 
$$2^R \left(\frac{2}{1+\tau\log N_{2}}\right)^{\lceil R/2 \rceil}\leq \frac{1}{2}\left(\frac{8}{1+\tau\log N_{2}}\right)^{\lceil R/2 \rceil}.$$
The conclusion follows after taking a union bound over $\bsh\in  \Acal_{d,\alpha,\bsgamma}(N_2)$.
\end{proof}


\begin{theorem}\label{thm:2}
With probability at least $1-\varepsilon_2$ for $\varepsilon_2$ given by \eqref{eqn:failprob2},
$$    \mathrm{err}(\Hcal_{d,\alpha,\bsgamma}, L_2,  A_{N,\bsz_{\{1{:}R\}},\Acal_{d,\alpha,\bsgamma}(N_2)}) \leq     C_{d, \alpha}  \frac{ (1+\log_2 N_2)^{\frac{d-1}{2}}}{N_2^{ \alpha}}, $$ where $C_{d, \alpha} > 0$ is a constant independent of $N$ and $N_2$ but dependent on $d$ and $\alpha$.
\end{theorem}

\begin{proof}
     For $f\in \Hcal_{d,\alpha,\bsgamma}$, we have
    \begin{equation}\label{eqn:L2errordecomp}
         \Vert A_{N,\bsz_{\{1{:}R\}},\Acal_{d,\alpha,\bsgamma}(N_2)}(f)-f\Vert^2_{L_2} = \sum_{\bsh\in \Acal_{d,\alpha,\bsgamma}(N_2)} \abs{\widehat{f}_{N,\bsz_{\{1{:}R\}}}(\bsh)-\widehat{f}({\bsh})}^2+\sum_{\bsh\in \ZZ^d\setminus\Acal_{d,\alpha,\bsgamma}(N_2)} \abs{\widehat{f}({\bsh})}^2.
    \end{equation}
    The second term on the right hand side can be bounded by
    \begin{equation*}
        \sum_{\bsh\in \ZZ^d\setminus\Acal_{d,\alpha,\bsgamma}(N_2)} \abs{\widehat{f}({\bsh})}^2\leq \frac{1}{N^{2\alpha}_2}\sum_{\bsh\in \ZZ^d\setminus\Acal_{d,\alpha,\bsgamma}(N_2)} \abs{\widehat{f}({\bsh})}^2r_{2\alpha,\bsgamma}(\bsh)\leq  \frac{\|f\|^2_{d,\alpha,\bsgamma}}{N^{2\alpha}_2}.
    \end{equation*}
    For the first term, Lemma~\ref{cor:whp2} together with Lemma~\ref{lem:medtosum} shows that with probability at least $1-\varepsilon_2$, 
        \begin{align*}
         & \sum_{\bsh\in \Acal_{d,\alpha,\bsgamma}(N_2)} \abs{\widehat{f}_{N,\bsz_{\{1{:}R\}}}(\bsh)-\widehat{f}({\bsh})}^2 \\
         & \leq 2\sum_{\bsh\in \Acal_{d,\alpha,\bsgamma}(N_2)}  \mathrm{median}_{r\in\{1{:}R\}} \abs{\widehat{f}_{N,\bsz_r}(\bsh)-\widehat{f}({\bsh})}^2 \\
         & \leq 2 \sum_{r\in\mathcal{R}} \sum_{\bsh\in K_{d,\alpha,\bsgamma}(\bsz_r)} \abs{\widehat{f}_{N,\bsz_r}(\bsh)-\widehat{f}({\bsh})}^2\\
         & \leq 2  \sum_{r\in\mathcal{R}} \sum_{\bsh\in K_{d,\alpha,\bsgamma}(\bsz_r)} \left(\sum_{\bsl\in P^\perp_{N,\bsz_r}\setminus\{\bszero\} } \abs{\widehat{f}({\bsh+\bsl})}\right)^2 \\
         & \leq 2  \sum_{r\in\mathcal{R}} \sum_{\bsh\in K_{d,\alpha,\bsgamma}(\bsz_r)} \left(\sum_{\bsl\in P^\perp_{N,\bsz_r}\setminus\{\bszero\} } \abs{\widehat{f}({\bsh+\bsl})}^2 r_{2\alpha,\bsgamma}(\bsh+\bsl) \right) \left(\sum_{\bsl\in P^\perp_{N,\bsz_r}\setminus\{\bszero\} }\frac{1}{r_{2\alpha,\bsgamma}(\bsh+\bsl)} \right)\\
         & \leq 2 R \|f\|^2_{d,\alpha,\bsgamma}  \sup_{r \in \mathcal{R} }\sup_{\bsh\in K_{d,\alpha,\bsgamma}(\bsz_r)} \sum_{\bsl\in P^\perp_{N,\bsz_r}\setminus\{\bszero\} }\frac{1}{r_{2\alpha,\bsgamma}(\bsh+\bsl)}.
    \end{align*}
    Because $\bsh+\bsl\notin \Acal_{d,\alpha,\bsgamma}(N_2)$ for $\bsh\in K_{d,\alpha,\bsgamma}(\bsz_r)$ and $\bsl\in P^\perp_{N,\bsz_r}\setminus\{\bszero\}$,
 \begin{align*}
    \sup_{\bsh\in K_{d,\alpha,\bsgamma}(\bsz_r)} \sum_{\bsl\in P^\perp_{N,\bsz_r}\setminus\{\bszero\} }\frac{1}{r_{2\alpha,\bsgamma}(\bsh+\bsl)} & \leq \sup_{\bsh\in \ZZ^d} \sum_{\bsl'\in (\bsh+ P^\perp_{N,\bsz_r})\setminus \Acal_{d,\alpha,\bsgamma}(N_2)}\frac{1}{r_{2\alpha,\bsgamma}(\bsl')} \\
    & \leq  \sum_{m=0}^\infty \frac{1}{(2^m N_{2})^{2\alpha} } |(\bsh+ P^\perp_{N,\bsz_r})\cap \Acal_{d,\alpha,\bsgamma}(2^{m+1}N_{2})|.    
    \end{align*}
Letting $m_0 \in \mathbb{Z}$ be such that $2^{m_0-1} < N_2 \le 2^{m_0}$, we apply \eqref{fiber:bound} to each $r\in \mathcal{R}$ and bound
   \begin{align}\label{ineq_min}
       & \sup_{r\in \mathcal{R} } \sup_{\bsh \in K_{d,\alpha,\bsgamma}(\bsz_r)} \sum_{\bsl\in P^\perp_{N,\bsz_r}\setminus\{\bszero\} }\frac{1}{r_{2\alpha,\bsgamma}(\bsh+\bsl)} \nonumber \\
       & \leq \sup_{r \in \mathcal{R}} \sum_{m=0}^\infty \frac{1}{(2^m N_{2})^{2\alpha} }   \frac{2^{m_0+m}}{m_0+m} \min \left\{ \frac{(1 + (m_0+m) /2)^d}{\rho_{2 \alpha, \bsgamma}(\bm{z}_r )^{1/(2\alpha) }}, \prod_{j=1}^d (1 + \gamma_j^{1/(2\alpha) } (m_0+m) )  \right\} \nonumber \\ 
       & \leq \frac{2}{N_2^{2\alpha-1}} \sup_{r \in \mathcal{R}} \sum_{m=0}^\infty \frac{1}{ 2^{m (2\alpha-1)}  }   \frac{1}{m_0+m}  \frac{(1 + (m_0+m) /2)^d}{\rho_{2 \alpha, \bsgamma}(\bm{z}_r )^{1/(2\alpha) }}  \\
       & \leq \frac{2}{N_2^{2\alpha}} \sum_{m=0}^\infty \frac{1}{ 2^{m (2\alpha-1)}  }     \frac{(1 + (m_0+m) /2)^d}{ m_0 + m } \nonumber \\
       & \leq C'_{d, \alpha} \frac{ (1+\log_2 N_2)^{d-1}}{N_2^{2 \alpha }} \nonumber
 \end{align}
 for a positive constant $C'_{d, \alpha}$.
  Here we used $1\leq m_0 < 1 +  \log_2 N_2$ and
 \begin{align*}
 m_0^{-d+1} \sum_{m=0}^\infty \frac{1}{ 2^{m (2\alpha-1)}  } \frac{ ( 1 + (m_0+m)/2)^d}{m_0+m} & \leq \sum_{m=0}^\infty \frac{1}{ 2^{m (2\alpha-1)}  } \left(\frac{1}{m_0+m} + \frac{1}{2} \right) \left( \frac{1}{m_0} + \frac{1}{2}+ \frac{m}{2 m_0} \right)^{d-1} \\ 
 & \leq \frac{3^d}{2^d}\sum_{m=0}^\infty \frac{ (1 + m)^{d-1}}{ 2^{m (2\alpha-1)}  }   \\
 & \leq \frac{3^d(d-1)!}{2^d}  \sum_{m=0}^\infty \frac{1}{ 2^{m (2\alpha-1)}  }\binom{m+d-1}{d-1} \\ 
 & = \frac{ 3^d (d-1)! }{2^d(1- 2^{-2\alpha+1})^d}. 
 \end{align*}
 Our conclusion follows by putting the above bounds into \eqref{eqn:L2errordecomp}.
\end{proof}

\begin{remark}
If we set the number of repetitions $R$ to be proportional to $\log N$, the failure probabilities $\epsilon_1$ and $\epsilon_2$ in Theorem~\ref{thm:infty} and Theorem~\ref{thm:2} vanish at the rate
\[ \mathcal{O}\left( \frac{N}{(\log N)^{\log N}}\right).\]
This decay rate is super-polynomial; that is, it decays faster than any fixed power of $N^{-1}$. Consequently, the total number of function evaluations, which is of order $N\log N$, is asymptotically sufficient for the median lattice algorithm to achieve high reliability.

Furthermore, the median lattice algorithm can also be viewed as a variant of multiple lattice-based algorithms, as the generating vectors $\bsz_r$ for which the median aggregation is attained can be different for each frequency $\bsh\in \Acal_{d,\alpha,\bsgamma}(N_2)$. While similar probabilistic strategies employing multiple lattices have been proposed in \cite{CG26,K19, KV19}, the median lattice approach offers a distinct advantage: it eliminates the need for an explicit verification step to detect aliasing terms across all frequencies in $\mathcal{A}_{d,\alpha,\boldsymbol{\gamma}}(N_2)$. Instead, the median operation ensures that the approximations of all Fourier coefficients are simultaneously free from aliasing issues with high probability, simplifying the implementation and reducing the computational overhead.
\end{remark}

 
\begin{remark}
In \eqref{ineq_min}, we ignored the second expression in the minimum in the line above. Using this second term instead yields an $L_2$ approximation error bound of the same form as the $L_\infty$ bound obtained in Theorem~\ref{thm:infty}. Since the $L_2$ approximation error is bounded above by the $L_\infty$ approximation error, this approach does not lead to a new result.
\end{remark}

\subsection{$L_p$-approximation}\label{sec_Lp}

We prove Theorem~\ref{thm:main} in this section. Since the $L_p$-approximation error for $1 \le p \le 2$ is bounded by the $L_2$-approximation error and we can estimate $(1+\log_2 N_2)^{d-1}$ by $C_{d,\beta} N^{\beta}$ for any $\beta > 0$, this part follows from Theorem~\ref{thm:2}. The bound independent of the dimension for $p=\infty$ follows directly from Theorem~\ref{thm:infty}.

It remains to prove the result for $2 < p < \infty$. We use the inequality $\|g\|_{L_p} \le \|g\|_{L_2}^{2/p} \|g\|_{L_\infty}^{1-2/p}$, which holds for $g \in L_2 \cap L_\infty$. Since every function $f$ in the 
weighted periodic Korobov space \( \Hcal_{d,\alpha,\bsgamma}\) is bounded and continuous, we have $f \in L_\infty$. Further, $A(f) \in L_\infty$ and so the norm inequality can be applied. 

As above, we estimate $(1+\log_2 N_2)^{d-1} \le C_{d,\beta} N^\beta$ for any $\beta > 0$. From Theorem~\ref{thm:infty} we get the error bound $\|f - A(f) \|_{L_\infty} \le C N^{-\alpha + 1/2 + \beta}$, and from Theorem~\ref{thm:2} we get the error bound $\|f - A(f)\|_{L_2} \le C N^{-\alpha + \beta}$. Applying the norm interpolation inequality, we obtain the bound
\begin{equation*}
\|f- A(f)\|_{L_p} \le C N^{(-\alpha + \beta) 2/p} N^{(-\alpha + 1/2 + \beta) (1- 2/p)} %= C N^{-\alpha + \theta/2  + \beta} 
= C N^{-\alpha + \frac12 - \frac1p + \beta},
\end{equation*}
and hence
\begin{equation*}
\mathrm{err}(\Hcal_{d,\alpha,\bsgamma}, L_p, A) \le C N^{-\alpha + \frac12 - \frac1p + \beta}, \quad 2 \le p \le \infty.
\end{equation*}
Thus, the result follows.


\section*{Acknowledgements}

We, the authors of this paper, have all been greatly inspired by the work of Henryk Wo\'{z}niakowski, and some of us have had the privilege to cooperate with him. We would like to thank Henryk for his numerous contributions to Information-Based Complexity and related fields, for his support, and---what is most important---his friendship. 

The fourth author acknowledges the support of the Japan Society for the Promotion of Science (JSPS) KAKENHI Grant Number 23K03210. The fifth author acknowledges the support of the Austrian Science Fund (FWF) Project  P 34808/Grant DOI: 10.55776/P34808. For open access purposes, the authors have applied a CC BY public copyright license to any author accepted manuscript version arising from this submission.

\begin{thebibliography}{99}

 \bibitem{BGKS24}
Felix Bartel, Alexander~D. Gilbert, Frances~Y. Kuo, and Ian~H. Sloan.
\newblock Minimal subsampled rank-1 lattices for multivariate approximation
  with optimal convergence rate.
\newblock {\em arXiv preprint arXiv:2506.07729}, 2025.

\bibitem{BKTV17}
Glenn Byrenheid, Lutz K\"{a}mmerer, Tino Ullrich, and Toni Volkmer.
\newblock Tight error bounds for rank-1 lattice sampling in spaces of hybrid
  mixed smoothness.
\newblock {\em Numer. Math.}, 136(4):993--1034, 2017.

\bibitem{CDG25}
Mou Cai, Josef Dick, and Takashi Goda.
\newblock A lattice algorithm with multiple shifts for function approximation
  in {K}orobov spaces.
\newblock {\em arXiv:2511.09071}, 2025.

\bibitem{CG26}
Mou Cai and Takashi Goda.
\newblock A note on approximation in weighted {K}orobov spaces via multiple
  rank-1 lattices.
\newblock {\em arXiv preprint arXiv:2601.20290}, 2026.

\bibitem{CKNS20}
Ronald Cools, Frances~Y. Kuo, Dirk Nuyens, and Ian~H. Sloan.
\newblock Lattice algorithms for multivariate approximation in periodic spaces
  with general weight parameters.
\newblock In {\em 75 years of {M}athematics of {C}omputation}, volume 754 of
  {\em Contemp. Math.}, pages 93--113. Amer. Math. Soc., RI, 2020.

\bibitem{CKNS21}
Ronald Cools, Frances~Y. Kuo, Dirk Nuyens, and Ian~H. Sloan.
\newblock Fast component-by-component construction of lattice algorithms for
  multivariate approximation with {POD} and {SPOD} weights.
\newblock {\em Math. Comp.}, 90(328):787--812, 2021.

\bibitem{D04}
Josef Dick.
\newblock On the convergence rate of the component-by-component construction of
  good lattice rules.
\newblock {\em J. Complexity}, 20(4):493--522, 2004.

\bibitem{DG21}
Josef Dick and Takashi Goda.
\newblock Stability of lattice rules and polynomial lattice rules constructed
  by the component-by-component algorithm.
\newblock {\em J. Comput. Appl. Math.}, 382:Paper No. 113062, 16, 2021.

\bibitem{DGS22}
Josef Dick, Takashi Goda, and Kosuke Suzuki.
\newblock Component-by-component construction of randomized rank-1 lattice
  rules achieving almost the optimal randomized error rate.
\newblock {\em Math. Comp.}, 91(338):2771--2801, 2022.

\bibitem{DKP22}
Josef Dick, Peter Kritzer, and Friedrich Pillichshammer.
\newblock {\em Lattice {R}ules---{N}umerical {I}ntegration, {A}pproximation,
  and {D}iscrepancy}, volume~58 of {\em Springer Series in Computational
  Mathematics}.
\newblock Springer, Cham, 2022.

\bibitem{DKS13}
Josef Dick, Frances~Y. Kuo, and Ian~H. Sloan.
\newblock High-dimensional integration: the quasi-{M}onte {C}arlo way.
\newblock {\em Acta Numer.}, 22:133--288, 2013.

\bibitem{DSWW06}
Josef Dick, Ian~H. Sloan, Xiaoqun Wang, and Henryk Wo\'zniakowski.
\newblock Good lattice rules in weighted {K}orobov spaces with general weights.
\newblock {\em Numer. Math.}, 103(1):63--97, 2006.

\bibitem{G26}
Takashi Goda.
\newblock A randomized lattice rule without component-by-component
  construction.
\newblock {\em Math. Comp.}, 95(357):339--361, 2026.

\bibitem{GK26}
Takashi Goda and David Krieg.
\newblock A simple universal algorithm for high-dimensional integration.
\newblock {\em Numer. Math.}, 158(1):229--248, 2026.

\bibitem{GL22}
Takashi Goda and Pierre L'Ecuyer.
\newblock Construction-free median quasi--{M}onte {C}arlo rules for function
  spaces with unspecified smoothness and general weights.
\newblock {\em SIAM J. Sci. Comput.}, 44(4):A2765--A2788, 2022.

\bibitem{GSM24}
Takashi Goda, Kosuke Suzuki, and Makoto Matsumoto.
\newblock A universal median quasi--{M}onte {C}arlo integration.
\newblock {\em SIAM J. Numer. Anal.}, 62(1):533--566, 2024.

\bibitem{HW63}
Loo~Keng Hua and Yuan Wang.
\newblock {\em Applications of {N}umber {T}heory to {N}umerical {A}nalysis}.
\newblock Springer-Verlag, Berlin-New York; Kexue Chubanshe (Science Press),
  Beijing, 1981.
\newblock Translated from the Chinese.

\bibitem{K19}
Lutz K\"{a}mmerer.
\newblock Constructing spatial discretizations for sparse multivariate
  trigonometric polynomials that allow for a fast discrete {F}ourier transform.
\newblock {\em Appl. Comput. Harmon. Anal.}, 47(3):702--729, 2019.

\bibitem{KPV15}
Lutz K\"{a}mmerer, Daniel Potts, and Toni Volkmer.
\newblock Approximation of multivariate periodic functions by trigonometric
  polynomials based on sampling along rank-1 lattice with generating vector of
  {K}orobov form.
\newblock {\em J. Complexity}, 31(3):424--456, 2015.

\bibitem{KV19}
Lutz K\"{a}mmerer and Toni Volkmer.
\newblock Approximation of multivariate periodic functions based on sampling
  along multiple rank-1 lattices.
\newblock {\em J. Approx. Theory}, 246:1--27, 2019.

\bibitem{KPUU25}
David Krieg, Kateryna Pozharska, Mario Ullrich, and Tino Ullrich.
\newblock Sampling recovery in {$L_2$} and other norms.
\newblock {\em Math. Comp.}, Published Online, 2025.

\bibitem{KKNU19}
Peter Kritzer, Frances~Y. Kuo, Dirk Nuyens, and Mario Ullrich.
\newblock Lattice rules with random {$n$} achieve nearly the optimal
  {$\mathcal{O}(n^{-\alpha-1/2})$} error independently of the dimension.
\newblock {\em J. Approx. Theory}, 240:96--113, 2019.

\bibitem{KNR19}
Robert~J. Kunsch, Erich Novak, and Daniel Rudolf.
\newblock Solvable integration problems and optimal sample size selection.
\newblock {\em J. Complexity}, 53:40--67, 2019.

\bibitem{K03}
F.~Y. Kuo.
\newblock Component-by-component constructions achieve the optimal rate of
  convergence for multivariate integration in weighted {K}orobov and {S}obolev
  spaces.
\newblock {\em J. Complexity}, 19(3):301--320, 2003.

\bibitem{KMNN21}
Frances~Y. Kuo, Giovanni Migliorati, Fabio Nobile, and Dirk Nuyens.
\newblock Function integration, reconstruction and approximation using rank-1
  lattices.
\newblock {\em Math. Comp.}, 90(330):1861--1897, 2021.

\bibitem{KNW23}
Frances~Y. Kuo, Dirk Nuyens, and Laurence Wilkes.
\newblock Random-prime-fixed-vector randomised lattice-based algorithm for
  high-dimensional integration.
\newblock {\em J. Complexity}, 79:Paper No. 101785, 28, 2023.

\bibitem{KSW06}
Frances~Y. Kuo, Ian~H. Sloan, and Henryk Wo\'zniakowski.
\newblock Lattice rules for multivariate approximation in the worst case
  setting.
\newblock In {\em Monte {C}arlo and quasi-{M}onte {C}arlo methods 2004}, pages
  289--330. Springer, Berlin, 2006.

\bibitem{LH03}
Dong Li and Fred~J. Hickernell.
\newblock Trigonometric spectral collocation methods on lattices.
\newblock In {\em Recent {A}dvances in {S}cientific {C}omputing and {P}artial
  {D}ifferential {E}quations ({H}ong {K}ong, 2002)}, volume 330 of {\em
  Contemp. Math.}, pages 121--132. Amer. Math. Soc., Providence, RI, 2003.

\bibitem{N92}
Harald Niederreiter.
\newblock {\em Random {N}umber {G}eneration and {Q}uasi-{M}onte {C}arlo
  {M}ethods}, volume~63 of {\em CBMS-NSF Regional Conference Series in Applied
  Mathematics}.
\newblock Society for Industrial and Applied Mathematics (SIAM), Philadelphia,
  PA, 1992.

\bibitem{NP09}
Wojciech Niemiro and Piotr Pokarowski.
\newblock Fixed precision {MCMC} estimation by median of products of averages.
\newblock {\em J. Appl. Probab.}, 46(2):309--329, 2009.

\bibitem{NW08}
Erich Novak and Henryk Wo\'zniakowski.
\newblock {\em Tractability of {M}ultivariate {P}roblems. {V}ol. 1: {L}inear
  {I}nformation}, volume~6 of {\em EMS Tracts in Mathematics}.
\newblock European Mathematical Society (EMS), Z\"urich, 2008.

\bibitem{NC06}
Dirk Nuyens and Ronald Cools.
\newblock Fast algorithms for component-by-component construction of rank-1
  lattice rules in shift-invariant reproducing kernel {H}ilbert spaces.
\newblock {\em Math. Comp.}, 75(254):903--920, 2006.

\bibitem{P26}
Zexin Pan.
\newblock Automatic optimal-rate convergence of randomized nets using
  median-of-means.
\newblock {\em Math. Comp.}, 95(359):1415--1446, 2026.

\bibitem{PGK25}
Zexin Pan, Takashi Goda, and Peter Kritzer.
\newblock Universal {$L_2$}-approximation using median lattice algorithms.
\newblock {\em arXiv preprint arXiv:2509.24582}, 2025.

\bibitem{PKG25}
Zexin Pan, Peter Kritzer, and Takashi Goda.
\newblock {$L_2$}-approximation using median lattice algorithms.
\newblock {\em arXiv preprint arXiv:2501.15331}, 2025.

\bibitem{PO23}
Zexin Pan and Art~B. Owen.
\newblock Super-polynomial accuracy of one dimensional randomized nets using
  the median of means.
\newblock {\em Math. Comp.}, 92(340):805--837, 2023.

\bibitem{PO24}
Zexin Pan and Art~B. Owen.
\newblock Super-polynomial accuracy of multidimensional randomized nets using
  the median-of-means.
\newblock {\em Math. Comp.}, 93(349):2265--2289, 2024.

\bibitem{SJ94}
I.~H. Sloan and S.~Joe.
\newblock {\em Lattice methods for multiple integration}.
\newblock Oxford Science Publications. The Clarendon Press, Oxford University
  Press, New York, 1994.

\bibitem{SR02}
I.~H. Sloan and A.~V. Reztsov.
\newblock Component-by-component construction of good lattice rules.
\newblock {\em Math. Comp.}, 71(237):263--273, 2002.

\bibitem{SW98}
Ian~H. Sloan and Henryk Wo\'zniakowski.
\newblock When are quasi-{M}onte {C}arlo algorithms efficient for
  high-dimensional integrals?
\newblock {\em J. Complexity}, 14(1):1--33, 1998.

\end{thebibliography}






\end{document}
